% วิธีดำเนินการวิจัย ภาษาไทย ใช้ใน main_short_th.tex (ฉบับอังกฤษอยู่ใน 03_methodology.tex); แก้ให้ตรงกันทั้งสองฉบับ.
% ไม่ใช้ \emph เพราะฟอนต์ไทยไม่มีตัวเอียง.
\section{วิธีดำเนินการวิจัย}\label{sec:method}

งานวิจัยนี้เปรียบเทียบแหล่งความลึกภายในกระบวนการสร้างแบบจำลองห้องที่คงที่เพียงกระบวนการเดียว
(รูปที่~\ref{fig:pipeline}) สำหรับทุกห้องและทุกแหล่งความลึก กระบวนการจะ (1)~อ่านภาพสีและท่าทางกล้อง
(2)~สร้างแผนที่ความลึกของทุกเฟรมจากแหล่งความลึกนั้น (3)~แปลงความลึกเป็นเมตร หากแหล่งนั้นไม่ได้ให้ค่าเป็นเมตร
(4)~รวมทุกเฟรมเป็นแบบจำลองห้อง (mesh) และ (5)~วัดทั้งความลึกรายเฟรมและแบบจำลองห้องเทียบกับค่าอ้างอิงจาก Faro
ระหว่างแถวต่าง ๆ ของตารางผล มีเพียงขั้นที่ (2) และ (3) ที่เปลี่ยน โดยแต่ละขั้นเลือกได้ด้วยค่าตั้งค่าเพียงรายการเดียว
ส่วนขั้นตอนหลังจากนั้นใช้โค้ดเดียวกันและค่าตั้งค่าเดียวกันทุกครั้ง โมเดลที่ปรับจูนแล้วเข้าสู่กระบวนการที่ขั้น (2)
ในฐานะแหล่งความลึกแหล่งหนึ่งตามปกติ บทนี้อธิบายข้อมูล แต่ละขั้นตอน การปรับจูน ตัวชี้วัดและสถิติ และการทดลอง

\begin{figure*}[t]
\centering
\begin{tikzpicture}[font=\footnotesize, >=Stealth, node distance=5mm,
  box/.style={draw=black!65, rounded corners=2pt, align=center, text width=29mm, minimum height=12mm, inner sep=2pt},
  swap/.style={box, draw=blue!70!black, very thick, fill=blue!4},
  side/.style={box, dashed}]
\node[box]  (in)   {\textbf{เฟรม}\\ภาพสี และท่าทางกล้องจาก ARKit};
\node[swap, right=of in]  (src)  {\textbf{แหล่งความลึก}\\Faro, LiDAR ของ iPad หรือโมเดลจากภาพเดี่ยว};
\node[swap, right=of src] (aln)  {\textbf{การปรับสเกล}\\ไม่ปรับ, oracle รายเฟรม, ครั้งเดียวต่อห้อง หรือจุดเบาบาง};
\node[box, right=of aln]  (fus)  {\textbf{การรวมแบบ TSDF}\\voxel \SI{4}{cm} แล้วใช้ marching cubes และทำความสะอาด};
\node[box, right=of fus]  (ev3)  {\textbf{ประเมินแบบสามมิติ}\\Chamfer, F-score เทียบแบบจำลองจาก Faro};
\node[side, above=8mm of src] (ft) {\textbf{การปรับจูน}\\บนห้องฝึก ด้วยข้อมูลสอนจาก Faro หรือ LiDAR};
\node[box, below=of aln]  (ev2)  {\textbf{ประเมินรายเฟรม}\\AbsRel, \wone{} เทียบความลึกจาก Faro};
\draw[->] (in) -- (src); \draw[->] (src) -- (aln); \draw[->] (aln) -- (fus); \draw[->] (fus) -- (ev3);
\draw[->, dashed] (ft) -- node[right, align=left, font=\scriptsize] {ค่าน้ำหนัก\\ของโมเดล} (src);
\draw[->] (aln) -- (ev2);
\end{tikzpicture}
\caption{กระบวนการของงานวิจัย สองขั้นที่เน้นสีเป็นเพียงสองขั้นที่เปลี่ยนระหว่างแถวของตารางผล
โดยเปลี่ยนด้วยค่าตั้งค่ารายการเดียว ขั้นตอนหลังจากนั้นคงที่ทั้งหมด การปรับจูน (เส้นประ) ทำครั้งเดียวล่วงหน้า
และทำหน้าที่เพียงให้ค่าน้ำหนักของโมเดลจากภาพเดี่ยว}
\label{fig:pipeline}
\end{figure*}

\subsection{ข้อมูล ห้อง และค่าอ้างอิง}\label{sec:data}
ARKitScenes~\cite{arkitscenes} ให้ข้อมูลของห้องเดียวกันหลายอย่าง ได้แก่ วิดีโอสีจาก iPad~Pro
(ขนาด $640\times480$ พิกเซล) ความลึกจาก LiDAR ที่ ARKit สร้างขึ้น (ขนาด $256\times192$ พิกเซล
พร้อมค่าความเชื่อมั่นระดับต่ำ กลาง หรือสูงของทุกพิกเซล) ท่าทางกล้องของทุกเฟรมจาก visual-inertial odometry ของ ARKit
และความลึกจากเครื่องสแกนเลเซอร์ Faro~\cite{faro} ที่จัดให้ตรงกับท่าทางกล้องเหล่านั้น
(ขนาด $1920\times1440$ พิกเซล ประมาณสี่เฟรมต่อวินาที) \textbf{ค่าอ้างอิงคือความลึกจาก Faro ตามที่ชุดข้อมูลให้มา}
สำหรับการประเมินรายเฟรม ความลึกนี้ถูกลดขนาดเป็น $256\times192$ โดยใช้พิกเซลที่ใกล้ที่สุด
ค่าอ้างอิงทุกค่าจึงเป็นค่าที่วัดได้จริง สำหรับการประเมินแบบสามมิติ เนื่องจากชุดข้อมูลไม่มีแบบจำลองจากเลเซอร์ให้
เราจึงรวมความลึกจาก Faro ของทุกเฟรมด้วย voxel ขนาด 1~ซม. ครั้งเดียวต่อห้อง เป็นแบบจำลองอ้างอิง
ด้วยวิธีรวมในหัวข้อ~\ref{sec:fusion} เนื่องจากแบบจำลองอ้างอิงและแหล่งความลึกทุกแหล่งใช้ท่าทางกล้องชุดเดียวกัน
ความคลาดเคลื่อนของท่าทางกล้องจึงไม่อยู่ในผลการวัด และความต่างทุกอย่างระหว่างแถวมาจากแหล่งความลึกเท่านั้น
งานนี้ไม่ใช้ LiDAR ของ iPad เป็นค่าอ้างอิงเลย

\begin{table}[t]
\centering\small
\caption{ห้องทดสอบ ประเมินทั้งแบบสามมิติและแบบรายเฟรม จำนวนเฟรม คือจำนวนเฟรมของ Faro ที่ใช้ประเมินทุกแถว}
\label{tab:scenes}
\begin{adjustbox}{max width=\columnwidth}
\begin{tabular}{@{}llrr@{}}
\toprule
รหัสการสแกน & ห้อง & ขนาด (ม.) & จำนวนเฟรม \\
\midrule
42444474 & ครัวและห้องทานข้าว มีกระจก (ห้องที่ใช้ระหว่างพัฒนา) & $9.9\times5.1\times4.2$ & 236 \\
47333774 & ห้องนอนทั่วไป & $5.1\times4.9\times3.3$ & 375 \\
47115299 & ห้องนั่งเล่นขนาดใหญ่ & $8.1\times6.1\times2.5$ & 324 \\
42897743 & ห้องน้ำ (กระจกเงา กระเบื้อง) & $8.0\times3.4\times3.6$ & 448 \\
47429736 & ห้องนอนที่มีตู้เสื้อผ้า & $6.6\times6.4\times2.7$ & 504 \\
45261556 & ครัว (ตู้ผิวมันวาว) & $7.1\times6.1\times2.4$ & 465 \\
\bottomrule
\end{tabular}
\end{adjustbox}
\end{table}

ห้องทดสอบทั้งหกห้อง (ตารางที่~\ref{tab:scenes}) ครอบคลุมห้องนอนทั่วไป ห้องขนาดใหญ่
และห้องที่มีกระจกเงาหรือพื้นผิวมันวาว ห้อง 42444474 ถูกใช้ระหว่างการพัฒนากระบวนการ
ผลของห้องนี้จึงไม่นับเป็นข้อมูลที่กันไว้ ผลรายห้องจึงช่วยให้ผู้อ่านตรวจข้อสรุปทุกข้อกับอีกห้าห้องที่เหลือได้
ห้องสำหรับการปรับจูนสุ่มเลือก (seed~0) จากการสแกนใน ARKitScenes ที่ยาว 60 ถึง 120 วินาที
โดยตัดทุก visit ของห้องทดสอบออก visit คือการเก็บข้อมูลบ้านหลังหนึ่งในครั้งเดียว ซึ่งวิดีโอในครั้งนั้นแสดงห้องเดียวกัน
ได้ห้องฝึก 10 ห้องที่มีทั้งความลึกจาก Faro และ LiDAR อีก 14 ห้องที่มีเฉพาะ LiDAR
และ 3 ห้องจากกลุ่ม Validation ของชุดข้อมูลสำหรับเลือก checkpoint
การทดลองที่ 7 เพิ่มอีก \rsVN{} ห้องจากกลุ่ม Validation โดยเลือกหนึ่งวิดีโอต่อ visit สุ่มด้วย seed~0
และตัดทุก visit ของห้องฝึก ห้องเลือก checkpoint และห้องทดสอบออก ห้องเหล่านี้ไม่ได้ใช้ตัดสินใจเรื่องใดเลย
และประเมินแบบรายเฟรมเท่านั้น บนทุกเฟรมที่สองของ Faro (รวม \rsVFRAMES{} เฟรม)

\subsection{แหล่งความลึก}\label{sec:sources}
งานนี้เปรียบเทียบแหล่งความลึกสี่ประเภท (1)~ความลึกจาก Faro เอง ซึ่งแสดงว่ากระบวนการสูญเสียความแม่นยำเท่าใด
เมื่อความลึกสมบูรณ์ เพราะที่ voxel ขนาด 4~ซม. ผลยังต่างจากแบบจำลองอ้างอิงที่ 1~ซม.
และความต่างนี้คือความสูญเสียจากการแบ่งช่องของกระบวนการเอง (2)~ความลึกจาก LiDAR ของ iPad ตามที่ ARKit บันทึก
ซึ่งเป็นสิ่งที่อุปกรณ์รุ่น Pro ใช้อยู่ในปัจจุบัน (3)~โมเดลความลึกสัมพัทธ์ ซึ่งต้องปรับสเกลตามหัวข้อ~\ref{sec:align}
ได้แก่ \mbox{DA-v2} Large ในการทดลองหลัก และ \mbox{DA-v2} Small กับ MiDaS small~\cite{midas}
เมื่อทดลองเปลี่ยนขนาดโมเดล (4)~โมเดลความลึกแบบเมตริก ซึ่งใช้โดยไม่ปรับสเกล ได้แก่ \mbox{DA-v2} Metric-Indoor Large
แบบเดิม โมเดลที่ปรับจูนแล้วสามตัว (หัวข้อ~\ref{sec:ftsetup}) และ Depth Pro~\cite{depthpro}
ซึ่งได้รับความยาวโฟกัสจริงของกล้อง 533 พิกเซล อันเป็นเงื่อนไขที่ได้เปรียบที่สุดของโมเดลนี้
โค้ดที่เรียกใช้ Depth Pro ผ่านการตรวจกับห้องสังเคราะห์ที่รู้ความลึกแน่นอนแล้ว และบนห้าเฟรมของห้อง 47115299
ยังรันโดยให้โมเดลประมาณมุมมองภาพ (field of view) เอง แล้ววัดผลหลังตัดสเกลของแต่ละเฟรมออกด้วย
ARKitScenes เก็บภาพจาก iPad ที่ถือในแนวตั้งโดยหมุนไปหนึ่งในสี่รอบ โมเดลจากภาพเดี่ยวทุกตัวจึงได้รับภาพที่หมุนให้ตั้งตรงก่อน
แล้วจึงหมุนผลลัพธ์กลับ แหล่งความลึกทุกแหล่งถูกปรับขนาดด้วยพิกเซลที่ใกล้ที่สุดเป็นความละเอียดของ LiDAR คือ
$256\times192$ เพื่อให้ทุกแหล่งถูกรวมและวัดผลบนตารางพิกเซลเดียวกัน

\subsection{การปรับสเกล}\label{sec:align}
โมเดลความลึกสัมพัทธ์ให้ส่วนกลับของความลึกที่ยังไม่รู้สเกลและค่าเลื่อนของแต่ละภาพ (บทที่~\ref{sec:related})
การปรับสเกลคือการหาสเกลและค่าเลื่อน ที่ทำให้ส่วนกลับของความลึกจากโมเดล ตรงกับส่วนกลับของความลึกอ้างอิงชุดหนึ่งมากที่สุด
แล้วแปลงผลกลับเป็นความลึก การหาค่าบนส่วนกลับของความลึกแทนที่จะหาบนความลึกโดยตรง สอดคล้องกับวิธีที่โมเดลเหล่านี้ถูกฝึก
และในห้องที่ใช้ระหว่างพัฒนา การเปลี่ยนนี้ลด Chamfer distance ของ oracle รายเฟรมจาก 15.2 เหลือ 5.4~ซม.
การหาค่าใช้วิธีกำลังสองน้อยที่สุด (least squares) ทำซ้ำสามรอบ แต่ละรอบตัดพิกเซล 20\,\% ที่มีค่าคลาดเคลื่อนมากที่สุดออก
เพราะขอบของวัตถุทำให้เกิดค่าคลาดเคลื่อนขนาดใหญ่มากจำนวนหนึ่ง ซึ่งจะดึงผลการหาค่าให้เพี้ยน
วิธีปรับสเกลทั้งสี่แบบต่างกันเพียงความลึกอ้างอิงที่ใช้หาค่า (ตารางที่~\ref{tab:aligners})
แหล่งความลึกแบบเมตริกใช้แบบไม่ปรับเสมอ และกระบวนการจะไม่ยอมทำงานหากตั้งค่าเป็นอย่างอื่น

\begin{table}[t]
\centering\small
\caption{วิธีปรับสเกล ทุกวิธียกเว้นแบบไม่ปรับ หาสเกลหนึ่งค่าและค่าเลื่อนหนึ่งค่าบนส่วนกลับของความลึก
ด้วยวิธีกำลังสองน้อยที่สุดแบบตัดค่าผิดปกติ}
\label{tab:aligners}
\begin{tabularx}{\columnwidth}{@{}l>{\raggedright\arraybackslash}X@{}}
\toprule
วิธี & หาค่าจากข้อมูลใด และแถวนี้บอกอะไร \\
\midrule
ไม่ปรับ & ไม่หาค่าใดเลย ใช้กับ Faro, LiDAR และโมเดลแบบเมตริก \\
oracle รายเฟรม & ความลึกจาก Faro ของเฟรมนั้นเอง แสดงขีดสูงสุดที่รูปทรงจากโมเดลทำได้ (ใช้ในผลิตภัณฑ์จริงไม่ได้) \\
ครั้งเดียวต่อห้อง & ความลึกจาก Faro ของ 10 เฟรมที่กระจายตลอดการสแกน หาค่าครั้งเดียวแล้วใช้กับทุกเฟรม
เป็นขอบบนของการปรับเทียบเพียงครั้งเดียว \\
จุดเบาบาง & จุดที่รู้ความลึกเป็นเมตร 200 จุดต่อเฟรม ใช้แทนจุดเด่นของ VIO (บทที่~\ref{sec:related})
ARKitScenes ไม่ได้เก็บจุดเหล่านี้ไว้ จึงสุ่มจากพิกเซลของ LiDAR ที่มีความเชื่อมั่นสูง ซึ่งแม่นกว่าจุดที่คำนวณจากภาพ
จึงเป็นขอบบนของระบบที่ใช้ VIO \\
\bottomrule
\end{tabularx}
\end{table}

สำหรับวิธีจุดเบาบาง การหาค่าทำซ้ำสองรอบและตัดจุดที่แย่ที่สุด 10\,\% ออก
เฟรมที่มีจุดใช้งานได้น้อยกว่า 20 จุดจะใช้สเกลและค่าเลื่อนของเฟรมก่อนหน้า
เช่นเดียวกับที่ระบบจริงจะทำเมื่อการติดตามภาพหลุดไปชั่วครู่

\subsection{การรวมเป็นแบบจำลองห้อง}\label{sec:fusion}
แผนที่ความลึกถูกรวมด้วยวิธี TSDF (บทที่~\ref{sec:related}) ใน Open3D~\cite{open3d} โดยใช้ voxel ขนาด 4~ซม.
ระยะตัด (truncation) 12~ซม. หรือสาม voxel และใช้ความลึกตั้งแต่ 0.1 ถึง 5~ม.
ทุกแหล่งความลึกรวมเฉพาะพิกเซลที่ Faro มีความลึก เพื่อไม่ให้แหล่งใดเสียคะแนนจากการมองเห็นส่วนของห้องที่ค่าอ้างอิงไม่ครอบคลุม
และทุกพิกเซลมีน้ำหนักเท่ากัน (มีเพียงการทดลองที่ 5 ที่ถ่วงน้ำหนักพิกเซลตามความเชื่อมั่นของ LiDAR)
พื้นผิวถูกสกัดด้วย marching cubes~\cite{marchingcubes} แล้วลบชิ้นส่วนที่เชื่อมกันน้อยกว่า \num{1000} สามเหลี่ยมออก
เพราะเป็นสัญญาณรบกวนที่ลอยอยู่ ทุกแถวใช้ท่าทางกล้องจาก ARKit ที่มากับชุดข้อมูล

\subsection{การปรับจูน}\label{sec:ftsetup}
โมเดลที่ปรับจูนจาก \mbox{DA-v2} Metric-Indoor Large มีสามตัว ใช้สูตรการฝึกเดียวกัน และต่างกันเพียงข้อมูลสอน
\ftfaro{} เรียนจากความลึกของ Faro บนห้องฝึก 10 ห้อง (\num{3594} เฟรม) \ftlidar{} เรียนจากความลึกของ LiDAR ของ iPad
บน 10 ห้องเดียวกัน (\num{8656} เฟรม) และ \ftmain{} ซึ่งเป็นโมเดลหลัก เรียนจากความลึกของ LiDAR บนทั้ง 24 ห้อง
(\num{19896} เฟรม) พิกเซลของ LiDAR ที่มีความเชื่อมั่นต่ำถูกตัดออก และไม่ใช้ข้อมูลสอนที่ไกลเกิน 10~ม.
ครูทั้งสองแบบเห็นเฟรมไม่เหมือนกัน เพราะความลึกจาก Faro มีประมาณสี่เฟรมต่อวินาที ขณะที่ความลึกจาก LiDAR
มากับภาพสีทุกเฟรม ซึ่งเราเลือกใช้ทุกเฟรมที่สาม (ประมาณสิบเฟรมต่อวินาที) \ftlidar{} จึงเห็นมุมมองที่ต่างกัน
มากกว่า \ftfaro{} ประมาณ 2.4 เท่า

ส่วนเข้ารหัสภาพถูกตรึงไว้ และฝึกเฉพาะส่วนถอดรหัส DPT คือ neck และ head (บทที่~\ref{sec:related})
ภาพฝึกแต่ละภาพถูกหมุนให้ตั้งตรง ปรับขนาดให้ด้านสั้นยาว 518 พิกเซล แล้วตัดเป็น $518\times518$ ที่ตำแหน่งสุ่ม
โดยไม่ขยายภาพเลย เพราะการขยายภาพทำให้วัตถุดูใหญ่ขึ้นและจะสอนระยะเป็นเมตรที่ผิด การพลิกภาพแนวนอน
(ความน่าจะเป็น 0.5) และการสุ่มปรับสี (ความแรง 0.2) ช่วยเพิ่มความหลากหลาย ฟังก์ชันสูญเสียคือ SiLog
(บทที่~\ref{sec:related}) ซึ่งสำหรับแต่ละภาพ คือความแปรปรวนของผลต่างระหว่างลอการิทึมของความลึกที่ทำนายกับของข้อมูลสอน
ซึ่งวัดรูปทรง บวกกับ 0.15 เท่าของกำลังสองของค่าเฉลี่ยของผลต่างนั้น ซึ่งวัดสเกลโดยรวม
แล้วนำรากที่สองของผลรวมคูณด้วย 10 ตัวหาค่าเหมาะที่สุด (optimiser) คือ AdamW~\cite{adamw}
ซึ่งเป็นรูปแบบหนึ่งของ Adam ที่แยกการลดค่าน้ำหนัก (weight decay) ออกจากขั้นตอนตามเกรเดียนต์
ใช้อัตราการเรียนรู้ \num{5e-5} weight decay 0.01 เพิ่มอัตราการเรียนรู้แบบเส้นตรงในช่วง 200 ขั้นแรก
แล้วลดลงแบบโคไซน์จนเป็นศูนย์~\cite{sgdr} แต่ละขั้นใช้ภาพ 8 ภาพ (ครั้งละ 2 ภาพ และสะสมเกรเดียนต์ 4 ครั้ง)
และจำกัดขนาดของเกรเดียนต์ไม่เกิน 1 การฝึกใช้ \num{6000} ขั้นสำหรับ \ftfaro{} และ \ftlidar{}
และ \num{12000} ขั้นสำหรับ \ftmain{} ทุก 500 ขั้น โมเดลจะถูกวัด AbsRel เทียบความลึกจาก Faro
บนทุกเฟรมที่ห้าของห้องเลือก checkpoint ทั้งสามห้อง แล้วเก็บ checkpoint ที่มี AbsRel ต่ำที่สุด
ดังนั้นค่าน้ำหนักของโมเดลที่เรียนจาก LiDAR ได้มาจาก LiDAR เท่านั้น แต่การเลือก checkpoint ใช้ความลึกจาก Faro
ของสามห้องนี้ การฝึก \num{6000} ขั้นใช้เวลาประมาณหนึ่งชั่วโมงบน GPU รุ่น NVIDIA RTX~3070
ขณะทดสอบ โมเดลที่ปรับจูนแล้วถูกใช้เหมือนโมเดลเดิมทุกประการ คือรับภาพสีและให้ความลึกเป็นเมตร
ไม่ปรับสเกล และเตรียมภาพนำเข้าแบบเดียวกัน

\subsection{ตัวชี้วัดและสถิติ}\label{sec:metrics}
\textbf{แบบสามมิติ} สุ่มจุด \num{200000} จุดอย่างสม่ำเสมอบนแบบจำลองแต่ละอัน ความแม่นยำ (accuracy)
คือระยะเฉลี่ยจากจุดบนแบบจำลองที่สร้างได้ถึงจุดที่ใกล้ที่สุดบนแบบจำลองอ้างอิง ความครบถ้วน (completeness)
คือระยะเฉลี่ยในทิศทางกลับกัน และ Chamfer distance~\cite{chamfer} คือค่าเฉลี่ยของทั้งสอง
precision, recall และ F-score~\cite{tanks} (บทที่~\ref{sec:related}) คำนวณที่ระยะเกณฑ์ 2, 5 และ 10~ซม.
โดยใช้ \fscore{} เป็นตัวเลขหลัก สำหรับวัตถุประสงค์ข้อ (3) ถือว่ารองรับงานได้เมื่อ recall ที่ระยะเกณฑ์ของงานนั้นอย่างน้อย 0.9
รองรับได้บางส่วนเมื่ออยู่ระหว่าง 0.75 ถึง 0.9 และรองรับไม่ได้เมื่อต่ำกว่า 0.75
\textbf{แบบรายเฟรม} แผนที่ความลึกแต่ละภาพถูกเทียบกับความลึกจาก Faro บนตาราง $256\times192$
เฉพาะพิกเซลที่ความลึกจาก Faro อยู่ระหว่าง 0.1 ถึง 5~ม. และแหล่งความลึกให้ค่าออกมา
ค่าที่ทำนายนอกช่วงนี้จะถูกตัดให้อยู่ในช่วงแทนที่จะถูกทิ้ง โมเดลที่วางผนังไกลเกินไปจึงยังถูกวัดผลที่ผนังนั้น
งานนี้รายงาน AbsRel, \wone{} และรากที่สองของค่าเฉลี่ยกำลังสองของความคลาดเคลื่อน (RMSE) ของความลึก~\cite{eigen14}
โดยเฉลี่ยในแต่ละห้องก่อน แล้วจึงเฉลี่ยข้ามห้อง
\textbf{สถิติ} ตารางแสดงค่าเฉลี่ย $\pm$ ส่วนเบี่ยงเบนมาตรฐานข้ามห้อง เนื่องจากหน่วยที่เป็นอิสระต่อกันคือห้อง ไม่ใช่เฟรม
การเปรียบเทียบสองโมเดลจึงใช้ค่ารายห้องกับการทดสอบ Wilcoxon signed-rank แบบ exact สองทาง~\cite{wilcoxon}
การทดสอบนี้เรียงลำดับผลต่างรายห้องตามขนาด แล้วตรวจว่าลำดับของผลต่างที่เป็นบวกและลบสมดุลกันหรือไม่
คำว่า exact หมายถึงคำนวณค่า p จากทุกรูปแบบที่เป็นไปได้ของเครื่องหมาย ไม่ได้ใช้ค่าประมาณ
ซึ่งสำคัญเมื่อมีห้องน้อย เมื่อมีหกห้อง ค่า p ที่เล็กที่สุดที่เป็นไปได้คือ $2/64 \approx 0.031$
ความต่างจึงมีนัยสำคัญที่ระดับ 5\,\% ได้ก็ต่อเมื่อโมเดลหนึ่งชนะครบทั้งหกห้อง

\subsection{การทดลอง}\label{sec:experiments}
ตารางที่~\ref{tab:exps} แสดงการทดลองและวัตถุประสงค์ที่แต่ละการทดลองตอบ การทดลองที่ 6 และ 7 ตอบคำถามหลัก
(วัตถุประสงค์ข้อ 1 และ 2) การทดลองที่ 1 ระบุตำแหน่งของโมเดลเทียบกับแหล่งความลึกที่ใช้ได้ขณะใช้งาน
(วัตถุประสงค์ข้อ 3) และการทดลองที่ 2--5 ทดสอบว่ากระบวนการไวต่อค่าตั้งค่าของตัวเองเพียงใด
เวลาในการทดลองที่ 4 วัดด้วย PyTorch ที่ไม่ได้ปรับแต่งบนแล็ปท็อป Apple Silicon และใช้เปรียบเทียบเป็นอัตราส่วนเท่านั้น

\begin{table}[t]
\centering\small
\caption{การทดลอง ในแต่ละการทดลองเปลี่ยนเฉพาะตัวแปรที่ระบุ ส่วนอื่นคงที่ทั้งหมด}
\label{tab:exps}
\begin{tabularx}{\columnwidth}{@{}ll>{\raggedright\arraybackslash}X@{}}
\toprule
การทดลอง & ข้อ & ตัวแปร ห้อง และการประเมิน \\
\midrule
6 & 1, 2 & โมเดลเดิมเทียบกับ \ftmain{}; \ftfaro{}, \ftlidar{} และ Depth Pro; ห้องทดสอบ 6 ห้อง สามมิติและรายเฟรม \\
7 & 1 & โมเดลเดิมเทียบกับ \ftmain{} โดยมี LiDAR ของ iPad เป็นแถวอ้างอิง; อีก \rsVN{} ห้อง รายเฟรม \\
1 & 3 & แหล่งความลึกและวิธีปรับสเกล: Faro, LiDAR, \mbox{DA-v2} Large กับการปรับแต่ละแบบ, \mbox{DA-v2} Metric แบบเดิม;
ห้องทดสอบ 6 ห้อง สามมิติและรายเฟรม \\
-- & 4 & การวิเคราะห์สเกลรายเฟรม (หัวข้อ~\ref{sec:scaleanalysis}); ห้องทดสอบ 6 ห้อง ทุกเฟรมที่ห้า \\
2--5 & -- & ขนาด voxel; ระยะห่างของเฟรม; ขนาดโมเดล; ความเชื่อมั่นของ LiDAR; ห้องทดสอบ 6 ห้อง สามมิติ \\
\bottomrule
\end{tabularx}
\end{table}

\subsection{การวิเคราะห์สเกลรายเฟรม}\label{sec:scaleanalysis}
วัตถุประสงค์ข้อ (4) ถามว่าความคลาดเคลื่อนที่การปรับจูนแก้ไม่ได้คืออะไร และมาจากเวลาหรือจากมุมมอง
บนทุกเฟรมที่ห้าของห้องทดสอบทั้งหก (471 เฟรม) เราคำนวณสำหรับโมเดลแบบเมตริกเดิมและ \ftmain{}
อัตราส่วนของความลึกจาก Faro ต่อความลึกที่ทำนาย โดยใช้ค่ามัธยฐานบนพิกเซลของเฟรมนั้น
ค่ามัธยฐานของอัตราส่วนนี้บนเฟรมทั้งหมดของห้องบอกว่าสเกลของโมเดลคลาดไปเท่าใด
และอัตราส่วนระหว่างเปอร์เซ็นไทล์ที่ 95 กับที่ 5 ของอัตราส่วนนี้บนเฟรมเหล่านั้น บอกว่าสเกลที่ต้องการแกว่งภายในห้องเดียวมากเพียงใด
จากนั้นคำนวณ AbsRel เมื่อให้ความช่วยเหลือด้านสเกลเพิ่มขึ้นทีละระดับ ได้แก่ ไม่ช่วยเลย
สเกลและค่าเลื่อนหนึ่งชุดต่อห้องที่หาจาก 10 เฟรมตามตารางที่~\ref{tab:aligners}
สเกลที่ถูกต้องของทุกเฟรม (ค่ามัธยฐานของอัตราส่วนของเฟรมนั้น ไม่มีค่าเลื่อน) และสเกลกับค่าเลื่อนแบบ oracle รายเฟรม
สำหรับโมเดลความลึกสัมพัทธ์ \mbox{DA-v2} Large เราหาสเกลแบบ oracle รายเฟรมบนเฟรมชุดเดียวกัน
แล้วคำนวณสัมประสิทธิ์สหสัมพันธ์ของ Pearson~\cite{pearson} ซึ่งวัดความสัมพันธ์เชิงเส้นในช่วง $-1$ ถึง $+1$
ระหว่างสเกลนี้กับเวลาของเฟรมในการสแกน และกับค่ามัธยฐานของความลึกจาก Faro ในเฟรมนั้น
หากเป็นการสะสมความคลาดเคลื่อน (drift) ตามเวลา สหสัมพันธ์กับเวลาจะมีเครื่องหมายเดียวกันในทุกห้อง
หากขึ้นกับมุมมอง สหสัมพันธ์กับความลึกของฉากตรงหน้ากล้องจะมีทิศทางเดียวกันอย่างสม่ำเสมอ
